% !PRINTED 772
\weberpage{772}

%<<CONT
ganzzahliges $x$ in eine ganze Zahl übergeht, die, mit $p$ multiplicirt, nicht $= 1$ sein kann.

Hiermit ist die Irreducibilität von $f_n(x)$ allgemein bewiesen.

\begin{center}
\rule{4cm}{0.4pt}
\end{center}

Nachdem so die Irreducibilität der Kreistheilungsgleichung
$f_n(x) = 0$ allgemein nachgewiesen ist, folgt, wie im \S.\,134 des
ersten Bandes, dass sie auch irreducibel bleibt, wenn eine beliebige Einheitswurzel adjungirt wird, deren Grad zu $n$ relativ
prim ist.

\gesperrt{Kronecker} geht gewissermaassen den umgekehrten Weg\footnote{Mém. sur les facteurs irréd. de l'expression $(x^n-1)$. \gesperrt{Liouville}'s
Journal, Bd. 19 (1854). Einen auf der Theorie der höheren Congruenzen
beruhenden einfachen Beweis der Irreducibilität der allgemeinen Kreistheilungsgleichung hat \gesperrt{Dedekind} gegeben [\gesperrt{Crelle}'s Journal für Mathematik, Bd. 58 (1859)].}.

Er beweist zunächst unter der Voraussetzung, dass $n$ eine
Potenz einer Primzahl $p$ ist, die Irreducibilität von $f_n(x)$ nach
Adjunction einer Einheitswurzel, deren Grad nicht durch $p$
theilbar ist, und leitet daraus den allgemeinen Irreducibilitätsbeweis ab, wie wir jetzt noch kurz zeigen wollen.

\begin{center}
\rule{4cm}{0.4pt}
\end{center}

Es sei jetzt
\[
n = p^k
\]
eine Potenz einer Primzahl $p$ und
\weq{8}{f_n(x) = x^{p^{k-1}(p-1)} + x^{p^{k-1}(p-2)} + \cdots + x^{p^{k-1}} + 1}
die Function, deren Wurzeln die primitiven $n^{\mathrm{ten}}$ Einheitswurzeln
sind. Es sei ferner $\alpha$ irgend eine Einheitswurzel, deren Grad $m$
durch $p$ nicht theilbar ist, und $\nu$ der Grad des Körpers $R(\alpha)$.
\gesperrt{Es soll bewiesen werden, dass $f_n(x)$ im Körper $R(\alpha)$ irreducibel ist.}

Jede Zahl des Körpers $R(\alpha)$ lässt sich auf eine und nur
auf eine Weise in der Form darstellen:
\weq{9}{\varphi(\alpha) = m_0 + m_1 \alpha + \cdots + m_{\nu-1} \alpha^{\nu-1},}
worin $m_0, m_1, \dots, m_{\nu-1}$ rationale Zahlencoëfficienten sind.

Ist $\chi(x) = 0$ die rationale Gleichung niedrigsten Grades,
deren Wurzel $\alpha$ ist, und
\weq{10}{\chi(x) = x^\nu + c_1 x^{\nu-1} + \cdots + c_{\nu-1} x + c_\nu,}
%CONT>>
